\PassOptionsToPackage{sols}{handout}
\documentclass[main]{subfiles}
\setcounterrefbeforedot{chapter}{m-ws6-2}
\setcounterrefafterdot{section}{m-ws6-2}
\addtocounter{section}{-1}
\usepackage{solutions}

\begin{document}

\ifSubfilesClassLoaded{%
  \chaptermark{\chaptersix}
}{}

\section{More Integration by Substitution} \label{ws6-2}

\begin{problemonly}
    \begin{framed}

\textbf{Key Ideas}:

\begin{itemize}[topsep=0pt,partopsep=0pt,parsep=8pt,itemsep=0pt]

\item You should be able to use substitution to rewrite certain
  integrals in a more manageable form.  (For example, if we substitute
  $u = x + 2$ in the integral $\displaystyle \int x \sqrt{x + 2}\,dx$,
  then we can rewrite the integral as a difference of two simpler
  integrals.)

\item You should develop flexibility with integration by substitution
  and feel more comfortable recognizing when to use substitution.  For
  example, the following integrals look similar but are evaluated using
  very different steps:
  \[ \int \frac{1}{x+1} \, dx \qquad \int \frac{1}{x^2+1} \, dx \qquad
  \int \frac{x}{x^2+1} \, dx \qquad \int \frac{x^2 +1}{x} \, dx\]

\end{itemize}
\end{framed}
\end{problemonly}

\begin{enumerate}
\item 

  \note{The point here is to give students a sense of the types of
    integrals that we can do using just algebra (without substitution), so
    that we have a ``target'' for our later substitutions.}

  \textbf{Warm-Up.}  Evaluate the following integrals; you should
  \textit{not} need to use substitution, only algebra.

  \probonly{}
    \begin{enumerate}
    \item
      \begin{problem}[\stretch{0.5}]
        $\displaystyle \int (u - 5) \sqrt{u}\,du$
      \end{problem}

      \begin{solution} % originally by Henry Shin, modified by Janet
        If we multiply out the integrand, then we can do the integral
        easily: 
        \begin{align*}
          \int (u -5) \sqrt{u}\,du 
          & =  \int \left( u^{3/2} -5 u^{1/2} \right) du \\
          & =  \boxed{\frac{2}{5}u^{5/2} - \frac{10}{3} u^{3/2} + C} 
        \end{align*}
      \end{solution}

    \item
      \begin{problem}[\stretch{0.5}]
        $\displaystyle \int \frac{(t + 1)^2}{t}\,dt$
      \end{problem}

      \begin{solution}
        Let's first rewrite the numerator: 
        \begin{align*}
          \int \frac{(t + 1)^2}{t}\,dt &= \int \frac{t^2 + 2t + 1}{t}\,dt \\
          &= \int \left( \frac{t^2}{t} + \frac{2t}{t} + \frac{1}{t}
          \right) dt \\
          &= \int \left( t + 2 + \frac{1}{t} \right) dt \\
          &= \boxed{\frac{1}{2} t^2 + 2t + \ln |t| + C}
        \end{align*}
      \end{solution}

    \item
      \begin{problem}[\nstretch{0.5}]
        $\displaystyle \int \frac{2u + 3}{\sqrt{u}}\,du$
      \end{problem}

      \begin{solution}
        Here, we have a sum in the numerator, so we can split the
        integrand into a sum of two fractions:
        \begin{align*}
          \int \frac{2u + 3}{\sqrt{u}}\,du &= \int \left( \frac{2u}{\sqrt{u}}
            + \frac{3}{\sqrt{u}} \right) du \\
          &= \int \left( 2u^{1/2} + 3 u^{-1/2} \right) du \\
          &= \boxed{\frac{4}{3} u^{3/2} + 6 u^{1/2} + C}
        \end{align*}
      \end{solution}
    \end{enumerate}    
    \probonly{}

\item 

  \textit{So far, when we've used substitution, we've generally tried to
    pick $u$ so that we also see $du$ in the integral.  However, sometimes
    substitution is used simply to rewrite an integral in a friendlier
    form.  This is the case in each of the following problems.}

  \begin{enumerate}
  \item
    \begin{problem}[\vfill]
      $\displaystyle \int x \sqrt{x+5}\,dx$.  (Hint: Let $u = x + 5$.  What
      is $x$ in terms of $u$?)
    \end{problem}

    \begin{solution} % originally by Henry Shin, modified by Janet
      If $u = x+5$, then $x = u-5$ and $du =dx$, so
      \begin{align*}
        \int x \sqrt{x+5} \,dx & =  \int (u -5) \sqrt{u}\,du \\
        \intertext{This is just the indefinite integral from
          {{{{\#1}}(a)}}, so it's equal to} %
        & =  \frac{2}{5}u^{5/2} - \frac{10}{3} u^{3/2} + C \\
        & = \boxed{\frac{2}{5} (x+5)^{5/2} - \frac{10}{3} (x+5)^{3/2} + C}
      \end{align*}
    \end{solution}

  \item
    \begin{problem}[\vfill]
      $\displaystyle \int \frac{2x} {\sqrt{x - 3}}\,dx$
    \end{problem}

    \begin{solution} % originally by Henry Shin, modified by Janet
      If we substitute $u = x - 3$, then $x = u + 3$ and $du = dx$, so
      \begin{align*}
        \int \frac{2x}{\sqrt{x - 3}}\,dx & = \int \frac{2 (u +
          3)}{\sqrt{u}} du \\
          &= 2\int \frac{u +
          3}{\sqrt{u}} du \\
          &= 2\int (u^{1/2}+3u^{-1/2})\, du \\
        &= 2 \left( \frac{2}{3} u^{3/2} + 6 u^{1/2} \right) + C \\
        &=  \frac{4}{3} u^{3/2} + 12 u^{1/2} + C \\
        & = \boxed{\frac{4}{3}(x - 3)^{3/2} + 12 (x - 3)^{1/2} + C}
      \end{align*}
    \end{solution}

  \item
    \begin{problem}[\vfill\newpage]
      $\displaystyle \int \frac{7x}{x + 3}\,dx$
    \end{problem}

    \begin{solution}
      If we substitute $u = x + 3$, then $x = u - 3$ and $dx = du$, so
      \begin{align*}
        \int \frac{7x}{x + 3}\,dx &= \int \frac{7(u - 3)}{u}\,du \\
        & = 7\int \left( 1 - \frac{3}{u} \right) du \\
        & = 7(u - 3 \ln |u|) + C \\
        &= 7u - 21\ln|u| + C \\
        & = \boxed{7(x+3) - 21 \ln |x + 3| + C}
      \end{align*}
      (This can be written more simply as just $7x - 21 \ln |x + 3| + C$.) 
    \end{solution}

  \end{enumerate}

\item 

  \note{These are practice with both ``types'' of substitution (reversing
    the Chain Rule and simplifying algebraically).  They require quite a
    lot of flexibility.}

  Evaluate each of the following integrals.

  \begin{enumerate}

  \item
    \begin{problem}[\stretch{0.5}]
      $\displaystyle \int \frac{e^{\sqrt{x}}}{\sqrt{x}}\,dx$
    \end{problem}

    \begin{solution}
      It looks like a good substitution is \sub{$u = \sqrt{x} = x^{1/2}$}
      because it's the inner function in the composition $e^{\sqrt{x}}$,
      and we also see something related to $du = \frac{1}{2}
      x^{-1/2}\,dx$, namely \subd{$2\,du = x^{-1/2}\,dx =
        \dfrac{1}{\sqrt{x}}\,dx$}.  Then,
      \begin{align*}
        \int \frac{e^{\sub{\sqrt{x}}}}{\subd{\sqrt{x}}}\,\subd{dx} &= 
        \int e^\sub{u} \cdot \subd{2\,du} \\
        &= 2 \int e^u\,du \\
        &= 2 e^u + C \\
        &= \boxed{2 e^{\sqrt{x}} + C}
      \end{align*}
    \end{solution}

  \item
    \note{This is quite challenging.  Often their first guess is $u =
      x^2 + 1$, and from there you can get them to see that this works
      well for $\displaystyle \int \frac{2x}{x^2 + 1}\,dx$, which
      suggests splitting up the given integral.}

    \begin{problem}[\vfill]
      $\displaystyle \int_0^1 \frac{2x+3}{x^2+1}\,dx$
    \end{problem}

    \begin{solution}
      We can rewrite this integral as a sum of two integrals and work on
      the two integrals separately.\footnote{Is it obvious when you look
        at this problem that this is the right first step?  Probably not!
        But, as you practice more, you'll start to develop your intuition
        about how to tackle integrals like this.}
      \begin{equation} \label{eq:substitution-split-into-two-1}
 \int_0^1
        \frac{2x+3}{x^2+1}\,dx = \int_0^1 \frac{2x}{x^2+1}\,dx + \int_0^1
        \frac{3}{x^2+1}\,dx
      \end{equation}
      The second integral is one you should just recognize:
      \begin{align*}
        \int_0^1 \frac{3}{x^2+1}\,dx &= 3 \arctan x \Big|_0^1 \\
        &= 3 \cdot \frac{\pi}{4} - 3 \cdot 0 \\
        &= \frac{3 \pi}{4}
      \end{align*}
      The first integral is not one we just recognize, but the
      substitution \sub{$u = x^2 + 1$} looks useful, both because this is
      the denominator and because we see \subd{$du = 2x\,dx$} in the
      integral.  As $x$ goes from $0$ to $1$, $u = x^2 + 1$ goes from $0^2
      + 1 = 1$ to $1^2 + 1 = 2$.
      \begin{align*}
        \int_0^1 \frac{\subd{2x}}{\sub{x^2+1}}\,\subd{dx} &= \int_1^2
        \frac{1}{\sub{u}}\,\subd{du } \\
        &= \ln |u| \Big|_1^2 \\
        &= \ln 2 - \ln 1 \\
        &= \ln 2
      \end{align*}
      Putting this all together with equation
      \eqref{eq:substitution-split-into-two-1}, $\displaystyle \int_0^1
      \frac{2x+3}{x^2+1}\,dx = \boxed{\ln 2 + \frac{3 \pi}{4}}$.
    \end{solution}

  \item 
    \begin{problem}[\vfill\newpage]
      $\displaystyle \int y \cdot \sqrt[3]{2y+5}\,dy$
    \end{problem}

    \begin{solution}
      A reasonable substitution to try is \sub{$u = 2y + 5$}, since it's
      the inner function in the composition $\sqrt[3]{2y + 5}$.  With this
      substitution, $du = 2\,dy$, so \subd{$\frac{du}{2} = dy$}.  We also
      see a $y$ by itself in the integral, but we can express this in
      terms of $u$: the equation $u = 2y + 5$ can be rewritten as
      \sube{$\frac{u - 5}{2} = y$}: 
      \begin{align*}
        \int \sube{y} \sqrt[3]{\sub{2y + 5}}\,\subd{dy} &= \int
        \sube{\frac{u - 5}{2}} \sqrt[3]{\sub{u}} 
        \cdot \subd{\frac{1}{2}\,du} \\
        &= \frac{1}{4} \int (u - 5) \sqrt[3]{u}\,du \\
        &= \frac{1}{4} \int \left( u^{4/3} - 5 u^{1/3} \right) du \\
        &= \frac{1}{4} \left( \frac{3}{7} u^{7/3} - \frac{15}{4}
          u^{4/3} \right) + C \\
        &= \boxed{\frac{1}{4} \left( \frac{3}{7} (2y + 5)^{7/3} -
            \frac{15}{4} (2y + 5)^{4/3} \right) + C}
      \end{align*}
    \end{solution}

  \item
    \begin{problem}[\vfill]
      $\displaystyle \int e^{3x} \sqrt{e^x + 2}\,dx$
    \end{problem}

    \begin{solution}
      A reasonable substitution to try is \sub{$u = e^x + 2$}, since it's
      the inner function in the composition $\sqrt{e^x + 2}$.  Then,
      \subd{$du = e^x\,dx$}, which might not look promising at first since
      we don't see this in the integral.  However, if we rewrite $e^{3x}$
      as $e^{2x} e^x$, then we do see $du$ in the integral:
      \begin{align*}
        \int e^{3x} \sqrt{e^x + 2}\,dx &= \int e^{2x} \subd{e^x}
        \sqrt{\sub{e^x + 2}}\,\subd{dx} \\
        \intertext{However, we still have to figure out what to do with
          $e^{2x}$.  We can rewrite it as $(e^x)^2$ or $e^x \cdot e^x$,
          and the substitution $u = e^x + 2$ can be rewritten as \sube{$u
            - 2 = e^x$}:} %
        &= \int \left( \sube{e^x} \right)^2 \subd{e^x}
        \sqrt{\sub{e^x + 2}}\,\subd{dx} \\
        &= \int \left( \sube{u - 2} \right)^2 \sqrt{\sub{u}}\, \subd{du} \\
        &= \int (u^2 - 4u + 4) u^{1/2}\,du \\
        &= \int \left( u^{5/2} - 4u^{3/2} + 4u^{1/2} \right)\,du \\
        &= \frac{2}{7} u^{7/2} - \frac{8}{5} u^{5/2} + \frac{8}{3}
        u^{3/2} + C \\
        &= \boxed{\frac{2}{7} (e^x + 2)^{7/2} - \frac{8}{5} (e^x +
          2)^{5/2} + \frac{8}{3} (e^x + 2)^{3/2} + C}
      \end{align*}
    \end{solution}

  \item
    \begin{problem}[\vfill]
      $\displaystyle \int \frac{e^x}{e^{2x} + 2 e^x + 1}\,dx$
    \end{problem}

    \begin{solution}
      There's not an obvious substitution to use here, but we might try
      \sub{$u = e^x$} because everything in the integral can be written in
      terms of that, and we see \subd{$du = e^x\,dx$} in the integral:
      \begin{align*}
        \int \frac{e^x}{e^{2x} + 2 e^x + 1}\,dx &= \int
        \frac{\subd{e^x}}{(\sub{e^x})^2 + 2 \sub{e^x} + 1}\,\subd{dx} \\ 
        &= \int \frac{1}{u^2 + 2u + 1}\,du \\
        &= \int \frac{1}{(u + 1)^2}\,du \\
        \intertext{Now, we see a composition of functions, so we might try
          substituting for the inner function: if \sub{$v = u + 1$}, then
          \subd{$dv = du$}, and 
          so the previous integral can be rewritten as}
        &= \int \frac{1}{\sub{v}^2}\,\subd{dv} \\
        &= \int v^{-2}\,dv \\
        &= - v^{-1} + C \\
        &= - \frac{1}{v} + C \\
        &= - \frac{1}{u + 1} + C \\
        &= \boxed{- \frac{1}{e^x + 1} + C}
      \end{align*}
    \end{solution}

  \end{enumerate}
  
  \pspace{\newpage}


\item  It is 10 am, and five ants have found their way into a picnic
  basket.  Ants are notorious followers, so ants from all over the
  vicinity follow their five brethren into the basket. The culinary treat
  awaiting them is unsurpassed elsewhere, so once the ants find their way
  into the basket they choose not to leave.  Suppose the rate at which
  ants are climbing into the basket is well modeled by $100 e^{-0.2t}$
  ants per hour, where $t=0$ is the benchmark hour of 10 am.
  \begin{enumerate}
  \item
    \begin{problem}[\vfill]
      How many ants are in the basket at 1 pm?
    \end{problem}

    \begin{solution}
      Since there are 5 ants at 10 am (which is $t = 0$), the number of
      ants at 1 pm ($t = 3$) is
      \begin{align*}
        5 + \int_0^3 100 e^{-0.2t}\,dt &= 5 + \left( -500
          e^{-0.2t} \Big|_0^3 \right) \\
        &= \boxed{5 - 500 e^{-0.6} + 500}
      \end{align*}
    \end{solution}

  \item
    \begin{problem}[\vfill]
      How many ants are in the basket $x$ hours after 10 am?
    \end{problem}

    \begin{solution}
      Using exactly the same idea as in the previous part, the number of
      ants $x$ hours after 10 am is
      \begin{align*}
        5 + \int_0^x 100 e^{-0.2t}\,dt &= 5 + \left( -500
          e^{-0.2t} \Big|_0^x \right) \\
        &= \boxed{5 - 500 e^{-0.2 x} + 500}
      \end{align*}
    \end{solution}

  \end{enumerate}

\pspace{\newpage}

\item 

  Suppose $g$ is a mysterious function, and the only thing you know
  about it is that $\displaystyle \int_0^{12} g(x)\,dx = \pi$.  For each
  of the following integrals, decide whether you have enough information
  to evaluate the integral; if so, evaluate it.
  \probonly{}
    \begin{enumerate}
    \item
      \begin{problem}[\stretch{0.3}]
        $\displaystyle \int_0^{12} g(3x)\,dx$
      \end{problem}

      \begin{solution}
        There is \fbox{not enough information}.  (As you can see in the
        next part, substituting $u = 3x$ doesn't help.)  We can understand
        why graphically; suppose $g$ looks like this.  The only
        information we're given is that the green area is $\pi$:
        \begin{center}
          \begin{tikzpicture}[declare
            function={f(\x)=(\x<=4*pi)*(\x/1.2^\x+1.5)+
              (\x>4*pi)*(0.5*(5.542254215512511 +
              abs(sin(deg(0.5*\x)))));}]
            \begin{axis}[ymin=0, xtick={12}, ytick=\empty, width=2.5in,
              height=1.5in, clip=false]
              \addplot+[domain=0:12, plusarea] { f(x) } \closedcycle; %
              \addplot+[domain=0:12*pi, samples=121] { f(x) } node[right]
              {$y = g(x)$}; %
            \end{axis}
          \end{tikzpicture}
        \end{center}
        The graph of $g(3x)$ is stretched horizontally to be $\frac{1}{3}$
        as wide, so here's how we picture $\displaystyle \int_0^{12}
        g(3x)$:
        \begin{center}
          \begin{tikzpicture}[declare
            function={f(\x)=(\x<=4*pi)*(\x/1.2^\x+1.5)+
              (\x>4*pi)*(0.5*(5.542254215512511 +
              abs(sin(deg(0.5*\x)))));}]
            \begin{axis}[ymin=0, xtick={12}, ytick=\empty, width=2.5in,
              height=1.5in, clip=false]
              \addplot+[domain=0:12, fillregion] { f(3*x) }
              \closedcycle; %
              \addplot+[domain=0:12*pi, samples=181] { f(3*x) }
              node[right] {$y = g(3x)$}; %
            \end{axis}
          \end{tikzpicture}
        \end{center}
        But there's no way for us to find this area because we can't
        express it simply in terms of the green area above, which is the
        only area we know.
      \end{solution}

    \item
      \begin{problem}[\stretch{0.3}]
        $\displaystyle \int_0^4 g(3x)\,dx$
      \end{problem}

      \begin{solution}
        We can evaluate this by substituting \sub{$u = 3x$}; then, $du =
        3\,dx$, so \subd{$\frac{du}{3} = dx$}.  Since $x$ goes from $0$ to
        $4$, $u = 3x$ goes from $0$ to $12$; thus,
        \begin{align*}
          \int_0^{4} g(\sub{3x})\,\subd{dx} &= \int_0^{12} g
          (\sub{u})\,\subd{\frac{du}{3}} \\
          &= \frac{1}{3} \int_0^{12} g(u)\,du \\
          &= \boxed{\frac{\pi}{3}}
        \end{align*}
        This should also make sense graphically; the graph of $g(3x)$ is
        obtained by stretching the graph of $g(x)$ horizontally by a
        factor of $\frac{1}{3}$:
        \begin{center}
          \begin{tikzpicture}[declare
            function={f(\x)=(\x<=4*pi)*(\x/1.2^\x+1.5)+
              (\x>4*pi)*(0.5*(5.542254215512511 +
              abs(sin(deg(0.5*\x)))));}]
            \begin{axis}[ymin=0, xtick={12}, ytick=\empty, width=2.5in,
              height=1.5in, clip=false]
              \addplot+[domain=0:12, plusarea] { f(x) } \closedcycle; %
              \addplot+[domain=0:12*pi, samples=121] { f(x) } node[right]
              {$y = g(x)$}; %
            \end{axis}
          \end{tikzpicture}
          \hspace{0.5in}
          \begin{tikzpicture}[declare
            function={f(\x)=(\x<=4*pi)*(\x/1.2^\x+1.5)+
              (\x>4*pi)*(0.5*(5.542254215512511 +
              abs(sin(deg(0.5*\x)))));}]
            \begin{axis}[ymin=0, xtick={4}, ytick=\empty, width=2.5in,
              height=1.5in, clip=false]
              \addplot+[domain=0:4, fillregion] { f(3*x) } \closedcycle; %
              \addplot+[domain=0:12*pi, samples=181] { f(3*x) }
              node[right] {$y = g(3x)$}; %
            \end{axis}
          \end{tikzpicture}
        \end{center}
        Therefore, the area under $g(3x)$ from $x = 0$ to $x = 4$ should
        be $\frac{1}{3}$ as large as the area under $g(x)$ from $x = 0$ to
        $x = 12$.
      \end{solution}

    \item
      \begin{problem}[\stretch{0.3}]
        $\displaystyle \int_0^{144} g(x^2)\,dx$
      \end{problem}

      \begin{solution}
        There is \fbox{not enough information}.  (If you try substituting
        $u = x^2$, you instead end up with the integral in
        {{{{{{\#6}}(d)}} on {the worksheet ``Integration by Substitution''}}}.)
      \end{solution}

    \item
      \begin{problem}[\stretch{0.3}]
        $\displaystyle \int_0^{2 \sqrt{3}} x g(x^2)\,dx$
      \end{problem}

      \begin{solution}
        We can evaluate this using the substitution \sub{$u = x^2$}; then,
        $du = 2x\,dx$, so \subd{$\frac{du}{2} = x\,dx$}.  As $x$ goes from
        $0$ to $2 \sqrt{3}$, $u = x^2$ goes from $0^2 = 0$ to $\left(2
          \sqrt{3} \right)^2 = 12$:
        \begin{align*}
          \int_0^{2 \sqrt{3}} \subd{x} g(\sub{x^2})\,\subd{dx} &=
          \int_0^{12} g(\sub{u}) \subd{\frac{du}{2}} \\
          &= \frac{1}{2} \int_0^{12}
          g(u)\,du \\
          &= \boxed{\frac{\pi}{2}}
        \end{align*}
      \end{solution}

    \item
      \begin{problem}[\stretch{0.3}]
        $\displaystyle \int_0^{144} x g(x^2)\,dx$
      \end{problem}

      \begin{solution}
        There is \fbox{not enough information}.
      \end{solution}

    \end{enumerate}
    \probonly{}

\end{enumerate}

\end{document}